% Zdroj: PDF priklady-podzim-2024.pdf, fyzická str. 24-25. Značka: specializace UI-SU. Zařazeno: S3.
\section{Rozhodovací strom}
\szzotazka{podzim 2024}{31}{Rozhodovací strom}

\noindent\textbf{Zadání.}\par
Máme data se třemi atributy $A$, $B$ a $C$ a cílovou třídou $Y$ uvedená v tabulkách níže.

\begin{enumerate}
  \item Na trénovacích datech uvedených v tabulce popište algoritmus tvoření rozhodovacího stromu
    a strom postavte (plný, bez omezení na počet dat v listech a při dělení).
  \item Pro atribut vybraný do kořene popište výpočet kritéria, na základě kterého jste daný atribut
    vybrali. (Potřebujete-li logaritmy, stačí přibližně, některé aproximace jsou v tabulce.)
    \begin{center}
    \begin{tabular}{lcccccc}
    \toprule
    $x$ & $1/5$ & $1/3$ & $2/5$ & $3/5$ & $2/3$ & $4/5$ \\
    $\approx \log_2(x)$ & $-2{,}32$ & $-1{,}58$ & $-1{,}32$ & $-0{,}74$ & $-0{,}58$ & $-0{,}32$ \\
    \bottomrule
    \end{tabular}
    \end{center}
  \item Dále popište některou standardní metodu prořezávání. Pokud metoda využívá validační data,
    použijte validační data v tabulce.
\end{enumerate}

\begin{center}
\begin{minipage}{0.45\linewidth}
\centering
Trénovací data:\par\medskip
\begin{tabular}{ccccc}
\toprule
 & $A$ & $B$ & $C$ & $Y$ \\
\midrule
0 & 1 & 2 & 6 & n \\
1 & 1 & 2 & 6 & n \\
2 & 3 & 2 & 8 & y \\
3 & 3 & 4 & 6 & y \\
4 & 3 & 4 & 8 & y \\
5 & 3 & 4 & 8 & n \\
6 & 3 & 4 & 8 & n \\
7 & 3 & 4 & 6 & y \\
\bottomrule
\end{tabular}
\end{minipage}\hfill
\begin{minipage}{0.45\linewidth}
\centering
Validační data:\par\medskip
\begin{tabular}{ccccc}
\toprule
 & $A$ & $B$ & $C$ & $Y$ \\
\midrule
8  & 1 & 2 & 8 & n \\
9  & 3 & 4 & 6 & y \\
10 & 3 & 4 & 6 & y \\
11 & 3 & 4 & 8 & y \\
\bottomrule
\end{tabular}
\end{minipage}
\end{center}

\medskip
\noindent\textbf{Náčrt řešení.}\par
Pro výběr atributu se užívá entropie nebo gini index. Entropie pro dvouhodnotovou pravděpodobnost je
definovaná
\[
  H(a,b) = -\frac{a}{a+b}\log_2\!\Big(\frac{a}{a+b}\Big) - \frac{b}{a+b}\log_2\!\Big(\frac{b}{a+b}\Big)
\]
např.
\[
  H(2/6,\, 4/6) = -\tfrac{1}{3}\log_2\!\big(\tfrac{1}{3}\big) - \tfrac{2}{3}\log_2\!\big(\tfrac{2}{3}\big)
  \approx \frac{-1{,}58}{3} - \frac{2\cdot(-0{,}58)}{3} = \frac{2{,}74}{3} \approx 0{,}91.
\]
Entropický zisk je rozdíl entropie před dělením a vážené entropie po dělení; entropický zisk dělení
dle $A$ v kořeni je:
\[
  H(4,4) - \frac{2}{8}\cdot 0 - \frac{6}{8}H(2,4) \approx 1 - \frac{6}{8}\cdot 0{,}91 \approx 0{,}32.
\]
Postavený strom vyjde stejný na základě entropie i gini, viz obrázek. Většinou je atribut k dělení zřejmý,
např. dělení podle $B$ v kořeni by vedlo k $[2,1]$, $[2,3]$, což má nižší entropický zisk kvůli nenulové
entropii v obou listech a více špatně klasifikovaných příkladů.

\medskip
Prořezávání na základě validačních dat vyhodnotí, že je vhodné nahradit pravý uzel $C$ listem, protože
sníží validační chybu. Alternativou může být použití cost complexity pruning.

\begin{center}\includegraphics[width=\linewidth]{rozhodovaci-strom-entropie-gini.png}\end{center}
